\newpage
\section{Deep Learning Drift Mitigation: Literature Review}

In the field of inertial navigation, many approaches have been considered 
to address the problem of orientation drift. The classical solution has been
to use sensor fusion techniques like Kalman filtering. Zhang et
al. \fixme{cite 1} implemented a quaternion-based KF 
with vector selection. The Kalman observation step used these vectors 
to improve the prediction of the state and its covariance. It was noted that
it worked for short time periods of both magnetic and accelerometer disturbances,
but for long periods, gyroscopic drift still could not be compensated by due to the accumulation of errors.
To provide a stronger reference, Alatise et al. \fixme{cite 2} used an extended KF (EKF) to fuse sensor IMU measurements with
visual data. This resulted in large improvements in regards to orientation
estimates, however the dependence on visual data increases the complexity and
is not applicable in visually disturbed environments.   